\documentclass{article}
\usepackage{pathfinder-readable}

\title{Risk-Averse Learning in a Strategic Allocation Mechanism}

\begin{document}
\maketitle

\begin{abstract}
The pair joins a paper on payments and learning in a strategic resource-allocation game with a paper on
cooperative optimisation of conditional value at risk (CVaR). The thread asked whether the second paper's
sampled CVaR method could make the first paper's mechanism risk-averse while preserving its guarantees. The
peers found a fixed-batch example in which an exact surrogate equilibrium fails to implement the true risk
objective and fails true-risk participation, even though both objectives have strong curvature. They also
derived conditional bounds on this error, but found no feasible oracle or convergence proof for the
resulting strategic learning process. The verifier therefore left the thread at DRAFT.
\end{abstract}

\section{Introduction}

In Q, agents choose private resource allocations and submit price bids. A manager uses quadratic payments to
align the game's equilibrium with the allocation that maximises the sum of their expected valuations. Under
the paper's conditions, the equilibrium is unique, payments balance, agents can do at least as well as
opting out, and a stochastic proximal best-response procedure converges to that equilibrium \cite{Q}.

In P, agents cooperate over one shared decision. Each aims to limit the upper tail of a local random loss
using CVaR, a measure of severe losses. They estimate CVaR from sampled losses, perturb decisions to
estimate gradients, exchange decisions with neighbours, and take projected steps. The paper bounds sampling
and smoothing errors and studies the limiting decision when batch sizes stay fixed \cite{P}.

The proposed connexion was to value Q's allocations by the negative CVaR of each agent's loss and use P's
empirical estimator to learn them. This raises two separate questions. Can Q's payment argument still
support the desired risk objective? If so, does learning from finite batches lead to an equilibrium with the
same properties? The peers tested both questions, without treating P's cooperative update as an algorithm
for Q's strategic game.

\section{What was done}

The peers first compared the decision problems. Q gives each agent its own allocation and price bid, whereas
P has agents work on copies of one shared decision. They therefore treated P as a possible source of risk
estimates and error bounds, not as a ready-made best-response procedure for Q.

Next they defined an agent's risk valuation. Let \(x_n\) be agent \(n\)'s allocation, \(\xi_n\) its random
input, and \(\psi_n(x_n,\xi_n)\) its satisfaction. Write \(J_n=-\psi_n\) for its loss, and let \(\beta_n\)
be the fraction of the worst losses used by CVaR. The proposed valuation was
\[
R_n(x_n)=-\operatorname{CVaR}_{\beta_n}
  \bigl(J_n(x_n,\xi_n)\bigr).
\]
For a deterministic payment \(t_n\), CVaR's response to adding a fixed amount gives risk-adjusted utility
\(R_n(x_n)-t_n\). This makes the sum of local risk valuations a possible target for Q's payment
construction, subject to Q's remaining conditions.

The peers then checked curvature and found that Q's assumption on expected valuation does not carry over
automatically. They constructed a flat CVaR valuation from a strongly curved expected valuation. To repair
that step, they considered the stronger assumption that every sampled loss is uniformly strongly convex.
They also tested what a fixed empirical batch changes, checked whether P's perturbations are feasible on Q's
allocation sets, and separated the sum of agents' risks from the risk of their aggregate loss. Last, they
used P's value error estimate to derive conditional guarantees at an exact equilibrium of a surrogate game.

\section{Results}

\subsection{Curvature and payments}

The curvature failure is elementary. For an allocation \(x\in[0,1]\), take two equally likely losses, \(0\)
and \(x^2-1\). Their mean is strongly convex in \(x\). With a worst-half tail, CVaR is zero for every \(x\)
in the interval. The associated risk valuation is flat. Thus the strong concavity that Q assumes for
expected valuation cannot simply be reused for CVaR. This is the counterexample checked by Ada and Emmy in
the ledger.

There is a conditional repair. If every sampled loss \(J_n(\cdot,\xi_n)\) is \(m\)-strongly convex for the
same \(m>0\), then its CVaR is \(m\)-strongly convex, so \(R_n\) is \(m\)-strongly concave. Indeed,
subtracting \(m\|x_n\|^2/2\) from every sampled loss leaves a convex function, and CVaR moves by the same
deterministic amount. This does not settle the mechanism on its own: the risk valuation can have corners,
and Q's other payment and constraint conditions still need to hold. Ada's curvature and cash-translation
checks support this narrower claim.

\subsection{A fixed-batch failure}

The central example keeps strong curvature. Let \(x\in[0,1]\), let \(\xi\) be equally likely to be zero or one, and set
\[
J(x,\xi)=\frac{x^2}{2}-\frac{3x}{4}+\xi x,
\qquad \beta=\frac12.
\]
Here \(\beta\) is the worst-half tail fraction. True CVaR loss is \(x^2/2+x/4\), with its minimum at
\(x=0\). With one observation in each batch, expected empirical CVaR loss is instead \(x^2/2-x/4\), with its
minimum at \(x=1/4\). Both losses are strongly convex. Under Q's payment rule with a zero price bid and
slack capacity, the latter allocation is an exact equilibrium of the surrogate game in the peers'
construction. Its true risk utility is \(-3/32\), below the zero-allocation outside utility of zero. Emmy
gave the example and then clarified that the participation failure concerns true CVaR utility, not every
realised payoff. It is an equilibrium counterexample; it is not a proof that Q's learning algorithm
converges to that equilibrium.

\subsection{Conditional accuracy at a surrogate equilibrium}

The peers also bounded what a well-defined surrogate could guarantee. Suppose all perturbation queries are
valid, the absolute loss for agent \(n\) is at most \(U_n\), and its CVaR is \(L_n\)-Lipschitz. Let
\(\delta_n\) be the smoothing radius, \(s_n\) the batch size, and \(\widetilde R_n\) the expected smoothed
empirical risk valuation. P's smoothing and empirical-distribution estimates give the uniform value error
\[
\sup_{x_n}|\widetilde R_n(x_n)-R_n(x_n)|\leq e_n,
\qquad
e_n=L_n\delta_n+
\frac{U_n\sqrt{2\pi}}{\beta_n\sqrt{s_n}}.
\]
The supremum is over the allocation domain on which the queries and bounds apply. Ada and Emmy independently
checked the ensuing payoff argument. At an exact equilibrium of Q's surrogate game, replacing \(\widetilde
R_n\) by \(R_n\) changes each comparison between the chosen action and a unilateral deviation by at most
\(2e_n\), because the payment is the same in each comparison. Q's opt-out comparison then places true risk
utility at least \(-2e_n\) below \(R_n(0)\), the risk valuation of allocating nothing. If the mechanism
exactly implements the surrogate optimum under Q's conditions, its gap from the best feasible sum of true
local risk valuations is at most \(2\sum_n e_n\). The payment balance still holds at that exact surrogate
equilibrium under the same mechanism conditions. None of these statements proves a bound for an unfinished
learning run.

The peers also obtained a conditional distance estimate from P's gradient-error bound. If a feasible
smoothing oracle exists and the smoothed strategic game's equilibrium operator is strongly monotone with
constant \(\nu>0\), comparison of its two equilibrium inequalities bounds the distance between the
fixed-batch and true-smoothed equilibria by
\[
\frac{1}{\nu}\left(\sum_n
\frac{2\pi d_n^2U_n^2}{\beta_n^2\delta_n^2s_n}
\right)^{1/2},
\]
where \(d_n\) is agent \(n\)'s allocation dimension. This estimate needs the stated strategic-game
conditions; it supplies no convergence theorem for a learning algorithm.

\section{Objections and limits}

P's perturbation method assumes that the feasible set contains a ball centred at the origin. Q's allocations
are nonnegative, and its electric-vehicle example includes flow equalities, so that assumption can fail even
before strategic learning is considered. The peers left open a perturbation method within the feasible set
or a justified extension of the loss outside it. Their fixed-batch example identifies the surrogate
objective's error without claiming that P's one-point update runs directly on Q's set.

There is also a choice of welfare target. Deterministic payments with individual CVaR valuations target the
sum of local CVaRs. This can differ from CVaR of total loss. Ada checked a two-agent example with strongly
convex sample losses: the local-risk optimum is \((0,1)\), the aggregate-risk optimum is \((1/3,2/3)\), and
the latter has an aggregate-risk advantage of \(1/3+4q/9>0\), where \(q\) is a positive curvature parameter
below \(1/2\). Ada also derived a pooling identity using payments contingent on realised losses. That
identity needs verifiable losses and is not yet an incentive-compatible mechanism. The thread did not check
outside literature, so it supports no claim of novelty.

\section{Why the thread stopped}

The verifier kept the work at DRAFT because the fixed-batch example is a substantive result and the
conditional error bounds are supported, but neither establishes the promised risk-averse strategic learning
guarantee. The smallest step that would restart this line is a feasible CVaR oracle for Q's nonnegative and
affine allocation sets, followed by an inexact proximal best-response convergence argument that tracks batch
and smoothing errors. That would connect the equilibrium statements above to an outcome agents can actually
learn.

\bibliographystyle{plainurl}
\bibliography{references}

\end{document}
